\documentclass[11pt]{article}

\usepackage[margin=1in]{geometry}
\usepackage{amsmath,amssymb,amsthm,mathtools}
\usepackage{enumitem}
\usepackage{hyperref}

\newtheorem{theorem}{Theorem}
\newtheorem{proposition}[theorem]{Proposition}
\newtheorem{corollary}[theorem]{Corollary}
\theoremstyle{definition}
\newtheorem{definition}[theorem]{Definition}
\newtheorem{example}[theorem]{Example}
\theoremstyle{remark}
\newtheorem{remark}[theorem]{Remark}

\newcommand{\PB}[2]{\{#1,#2\}}

\title{Two Notions of \emph{Integral of Motion}:\\
Symmetry-Generated Conservation Laws versus\\
the Universal $2n$ Constants of the General Solution}
\author{}
\date{}

\begin{document}
\maketitle

\begin{abstract}
Physicists commonly say that a mechanical system does or does not \emph{have}
integrals of motion, depending on whether its Hamiltonian possesses continuous
symmetries, and that possessing enough independent, mutually commuting
integrals (Liouville integrability) is a rare and dynamically important
property. At the same time, elementary ODE theory --- underlying, for
instance, a well-known remark of Landau and Lifshitz --- shows that
\emph{every} Hamiltonian system with $n$ degrees of freedom possesses $2n$
(or $2n-1$, if autonomous) independent integrals of motion, with no
reference to any symmetry whatsoever. This note explains why the two
statements are complementary rather than contradictory. We review the
symmetry-based (Noether) origin of dynamically useful integrals and the
Liouville--Arnold notion of integrability; we prove, via an explicit
construction using the phase-space flow map, the folklore fact that $2n$
(or $2n-1$) integrals always exist; we show, using an elementary bound
from symplectic linear algebra, that at most $n$ integrals can ever be in
mutual involution, so the ``always available'' integrals essentially never
contribute to integrability; and we illustrate the discussion with the
harmonic oscillator, a resonant/non-resonant pair of oscillators, and the
chaotic H\'enon--Heiles system.
\end{abstract}

\section{Introduction}

Two statements about ``integrals of motion'' circulate in classical
mechanics, and at first sight they seem to be in tension.

\begin{enumerate}[label=(\Alph*)]
\item \textbf{The symmetry statement.} A conserved quantity is tied to a
continuous symmetry of the system (Noether's theorem): momentum is
conserved because space is homogeneous, angular momentum because space is
isotropic, energy because the dynamics is time-translation invariant. A
generic system need not have any of these symmetries, and correspondingly
need not have (beyond perhaps energy) any conserved quantity at all. The
strongest and most useful version of ``having enough integrals'' is
\emph{Liouville integrability}: $n$ independent, pairwise Poisson-commuting
integrals for a system with $n$ degrees of freedom. Most systems with
$n\ge 2$ do not have this property; they are non-integrable, generically
chaotic.
\item \textbf{The counting statement.} Regardless of any of this, a system
with $n$ degrees of freedom always possesses $2n$ independent integrals of
motion (or $2n-1$ time-independent ones, if the system is autonomous). This
is not a statement about symmetry at all; it is a direct consequence of the
existence and uniqueness theorem for ordinary differential equations. It
appears already in Landau and Lifshitz's \emph{Mechanics} \cite{landau},
where it is used to motivate the very definition of an integral of motion.
\end{enumerate}

The puzzle is: if every system automatically has $2n-1$ (or $2n$)
integrals, in what sense can a system ``lack'' integrals, or fail to be
integrable? The resolution is that (A) and (B) are answers to different
questions. Statement (B) is a bare existence theorem: it guarantees that
\emph{some} functions with the defining property of a constant of motion
exist, locally, near any given point of phase space and time. It says
nothing about whether those functions are smooth beyond a small
neighbourhood, whether they have a closed form, whether they can be written
down without already having solved the equations of motion, or whether they
are useful for anything. Statement (A) is about a much more restrictive and
much more useful class of integrals: quantities with a fixed, symmetry-determined
functional form that remains valid however the ``non-symmetric'' part of the
dynamics is changed, and --- in the strongest case --- quantities that are
in involution with one another, which is what actually allows the equations
of motion to be solved.

The rest of this note makes this precise. Section~\ref{sec:prelim} fixes
notation. Section~\ref{sec:symmetry} recalls Noether's theorem and
Liouville integrability. Section~\ref{sec:universal} proves the counting
statement (B) as a theorem, with an explicit construction. Section~\ref{sec:involution}
proves that no matter how the $2n$ (or $2n-1$) integrals of
Section~\ref{sec:universal} are chosen, at most $n$ of them can ever be in
involution --- which is the precise sense in which counting integrals is
not the same as having integrability. Section~\ref{sec:examples} works
through the harmonic oscillator, a pair of oscillators (resonant and
non-resonant), and the chaotic H\'enon--Heiles system, to show concretely
how ``always existing'' and ``practically useless'' can coexist. Section~\ref{sec:discussion}
discusses the reconciliation and its relation to the ``isolating
integral'' terminology used in galactic dynamics.

\section{Preliminaries}\label{sec:prelim}

Let $(q,p)=(q_1,\dots,q_n,p_1,\dots,p_n)$ be canonical (Darboux) coordinates
on a $2n$-dimensional phase space, with symplectic form
$\omega=\sum_{i=1}^n dq_i\wedge dp_i$, and let
$H=H(q,p,t)$ be a Hamiltonian, possibly explicitly time-dependent. Hamilton's
equations are
\begin{equation}
\dot q_i=\frac{\partial H}{\partial p_i},\qquad
\dot p_i=-\frac{\partial H}{\partial q_i},\qquad i=1,\dots,n.
\end{equation}
For functions $F,G$ on phase space, the Poisson bracket is
\begin{equation}
\PB{F}{G}=\sum_{i=1}^n\left(\frac{\partial F}{\partial q_i}\frac{\partial G}{\partial p_i}
-\frac{\partial F}{\partial p_i}\frac{\partial G}{\partial q_i}\right).
\end{equation}
The Hamiltonian vector field of a function $F(q,p)$ is defined by
$\iota_{X_F}\omega=dF$; in Darboux coordinates,
$X_F=\sum_i\left(\partial_{p_i}F\,\partial_{q_i}-\partial_{q_i}F\,\partial_{p_i}\right)$,
and a direct computation gives $\omega(X_F,X_G)=\PB{F}{G}$.

\begin{definition}[Integral of motion]
A function $F=F(q,p,t)$ is an \emph{integral of motion} (or \emph{constant
of motion}, or \emph{conserved quantity}) if
\begin{equation}
\frac{dF}{dt}=\frac{\partial F}{\partial t}+\PB{F}{H}=0
\end{equation}
identically, i.e.\ $F$ is constant along every solution of Hamilton's
equations. Two integrals $F_1,\dots,F_k$ are \emph{functionally independent}
at a point if $dF_1,\dots,dF_k$ are linearly independent there, and they are
\emph{in involution} if $\PB{F_i}{F_j}=0$ for all $i,j$.
\end{definition}

\section{Integrals from symmetry, and Liouville integrability}\label{sec:symmetry}

The physically substantive notion of an integral of motion is tied to
invariance of the dynamics under a continuous transformation.

\subsection{Noether's theorem}

In the Hamiltonian formalism, Noether's theorem takes an almost
tautological, and for that reason very clean, form: if $Q=Q(q,p)$ has no
explicit time dependence, then
\begin{equation}
\frac{dQ}{dt}=\PB{Q}{H},
\end{equation}
so $Q$ is conserved if and only if $\PB{Q}{H}=0$ --- and this is exactly
the condition that $H$ is invariant under the flow generated by $Q$ (equivalently,
by the antisymmetry of the bracket, that $Q$ is invariant under the flow
generated by $H$). Conservation of $Q$ and invariance of $H$ under the
one-parameter group generated by $Q$ are the same statement. (The
Lagrangian version of Noether's theorem, with an infinitesimal
transformation $q_i\to q_i+\epsilon\eta_i(q,t)$ leaving the action
invariant up to a boundary term, is equivalent to this after a Legendre
transform; we will not need its more elaborate bookkeeping here.)

The simplest instances are \emph{cyclic coordinates}: if $H$ does not
depend on $q_i$ at all, then $\PB{p_i}{H}=-\partial H/\partial q_i=0$, so
$p_i$ is conserved. Whether a given coordinate is cyclic is a contingent
fact about the specific $H$ under study, not a general feature of
mechanical systems:
\begin{itemize}
\item translation invariance of $H$ in a Cartesian coordinate $\Rightarrow$
the conjugate linear momentum is conserved;
\item rotation invariance of $H$ about an axis $\Rightarrow$ the conjugate
angular momentum is conserved;
\item invariance of $H$ under time translation (i.e.\ $H$ autonomous)
$\Rightarrow$ $H$ itself (the energy) is conserved.
\end{itemize}
A generic Hamiltonian --- e.g.\ three bodies interacting gravitationally
with no special symmetry, or a double pendulum --- has none of the spatial
symmetries above, and only energy is conserved by this mechanism. Whether
\emph{any} further integral exists is then a separate, dynamical question,
not something guaranteed by the structure of mechanics as such. This is
precisely statement (A) of the introduction.

\subsection{Liouville--Arnold integrability}

The strongest useful sense in which a system ``has enough'' integrals is
integrability in the sense of Liouville.

\begin{theorem}[Liouville--Arnold, see \cite{arnold}]
Suppose a Hamiltonian system with $n$ degrees of freedom has $n$
functionally independent integrals of motion $F_1=H,F_2,\dots,F_n$ that are
pairwise in involution, $\PB{F_i}{F_j}=0$. Then, under mild completeness
conditions, the common level sets $\{F_i=c_i\}$ are (unions of) Lagrangian
tori, motion on each torus is quasi-periodic, and the system can be
integrated by quadratures in suitable action--angle variables.
\end{theorem}

The content of this theorem is that having $n$ integrals \emph{in
involution} is what actually lets one solve the equations of motion; the
number $n$ is not arbitrary, and involution --- not mere existence --- is
the crucial extra structure. Systems that fail to be Liouville integrable
(the generic case for $n\ge 2$) typically exhibit chaotic dynamics.

\section{The universal integrals of the general solution}\label{sec:universal}

We now turn to statement (B). The following construction is exactly the
one behind the folklore fact recorded, e.g., in \cite{landau}: the general
solution of the equations of motion contains $2n$ arbitrary constants
(fixed by initial data), and \emph{inverting} the solution to express those
constants in terms of $(q,p,t)$ produces $2n$ integrals of motion, valid
for any system whatsoever.

\begin{theorem}\label{thm:universal}
Let $H\in C^k(U\times I)$, $k\ge 1$, where $U\subseteq\mathbb R^{2n}$ is
open and $I\subseteq\mathbb R$ is an open interval, and suppose that for
every $(q_0,p_0)\in U$ and $t_0\in I$ the solution of Hamilton's equations
with $(q(t_0),p(t_0))=(q_0,p_0)$ exists, is unique, and remains in $U$ for
all $t\in I$ (no finite-time blow-up on $I$). Fix a reference time
$t_0\in I$. Then there exist $2n$ functions $C_1,\dots,C_{2n}\in
C^k(U\times I)$ such that:
\begin{enumerate}[label=(\roman*)]
\item each $C_a$ is an integral of motion: $C_a(q(t),p(t),t)$ is constant
along every solution $(q(t),p(t))$;
\item for each fixed $t\in I$, the map $(q,p)\mapsto\big(C_1(q,p,t),\dots,C_{2n}(q,p,t)\big)$
is a $C^k$ diffeomorphism of $U$ onto its image; in particular the $C_a$
are functionally independent.
\end{enumerate}
\end{theorem}

\begin{proof}
By the standard theorem on existence, uniqueness, and $C^k$ dependence of
solutions of ODEs on initial data (see e.g.\ \cite{hartman}), the time-$t$
flow map
\begin{equation}
\Phi_{t,t_0}:U\to U,\qquad \Phi_{t,t_0}(q_0,p_0)=\big(q(t;q_0,p_0,t_0),\,p(t;q_0,p_0,t_0)\big)
\end{equation}
is, for each $t\in I$, a $C^k$ diffeomorphism of $U$ onto its image, jointly
$C^k$ in $(t,q_0,p_0)$. Uniqueness of solutions gives the group law
$\Phi_{t_2,t_1}\circ\Phi_{t_1,t_0}=\Phi_{t_2,t_0}$ and $\Phi_{t_0,t_0}=\mathrm{id}$,
so $\Phi_{t,t_0}^{-1}=\Phi_{t_0,t}$. Define
\begin{equation}
C(q,p,t):=\Phi_{t,t_0}^{-1}(q,p)=\Phi_{t_0,t}(q,p)\in\mathbb R^{2n},
\end{equation}
with components $C_1,\dots,C_{2n}$; this is $C^k$ jointly in $(q,p,t)$ by
the same smooth-dependence theorem.

(i) If $(q(t),p(t))$ solves Hamilton's equations, then
$(q(t),p(t))=\Phi_{t,t_0}(q(t_0),p(t_0))$ for all $t\in I$, so
\begin{equation}
C(q(t),p(t),t)=\Phi_{t,t_0}^{-1}\big(\Phi_{t,t_0}(q(t_0),p(t_0))\big)=(q(t_0),p(t_0)),
\end{equation}
which does not depend on $t$. Hence each $C_a$ is an integral of motion.

(ii) For fixed $t$, $(q,p)\mapsto C(q,p,t)=\Phi_{t,t_0}^{-1}(q,p)$ is by
construction the inverse of the diffeomorphism $\Phi_{t,t_0}$, hence is
itself a diffeomorphism, with everywhere-invertible Jacobian. In
particular $dC_1\wedge\cdots\wedge dC_{2n}\ne 0$.
\end{proof}

Nothing in this argument refers to any symmetry of $H$: it holds for a
generic, non-autonomous, non-integrable, even chaotic system, exactly as
freely as for the harmonic oscillator. It is, in an honest sense, a
bookkeeping device: $C(q,p,t)$ simply records \emph{which initial condition
the trajectory through $(q,p)$ at time $t$ came from}. Knowing all $2n$
components of $C$ as explicit functions is equivalent to already having
solved the equations of motion.

\begin{corollary}[Autonomous reduction]\label{cor:autonomous}
If $H=H(q,p)$ has no explicit time dependence, then $\Phi_{t,t_0}$ depends
on $t,t_0$ only through $\tau=t-t_0$. Away from equilibria ($X_H\ne 0$),
the flow-box (straightening) theorem for non-vanishing vector fields
provides local coordinates $(\tau,y_1,\dots,y_{2n-1})$ in which the flow is
simply $\tau\mapsto\tau+t$, $y\mapsto y$. Then $C_{2n}:=t-\tau(q,p)$ is the
one integral that carries explicit $t$-dependence, while
$C_1,\dots,C_{2n-1}:=y_1(q,p),\dots,y_{2n-1}(q,p)$ are $2n-1$ functionally
independent integrals with \emph{no} explicit time dependence. Thus an
autonomous system with $n$ degrees of freedom has $2n-1$ time-independent
integrals of motion (plus one further, explicitly time-dependent one),
matching the count in \cite{landau}.
\end{corollary}

\begin{remark}
The reduction in Corollary~\ref{cor:autonomous} is only \emph{local}: away
from equilibria it holds near any given point, but the resulting $C_a$ need
not extend to smooth, single-valued functions on all of $U$. Example~\ref{ex:pair}
below exhibits this concretely.
\end{remark}

\section{Counting is not enough: an involution bound}\label{sec:involution}

The $2n$ (or $2n-1$) integrals produced by Theorem~\ref{thm:universal} vastly
outnumber the $n$ integrals required for Liouville integrability. One might
wonder whether some subset of them could simply be selected to establish
integrability of \emph{any} system. The following elementary fact from
symplectic linear algebra shows this is never possible beyond a hard bound,
regardless of how the integrals are chosen.

\begin{proposition}\label{prop:bound}
On a $2n$-dimensional symplectic manifold, at most $n$ functionally
independent functions can be pairwise in involution in a neighbourhood of a
point where their differentials are linearly independent.
\end{proposition}

\begin{proof}
Let $F_1,\dots,F_k$ be functionally independent and pairwise in involution
near $x$, so $dF_1(x),\dots,dF_k(x)$ are linearly independent and
$\PB{F_i}{F_j}=0$ for all $i,j$. Let $X_i:=X_{F_i}$; since $\omega$ is
non-degenerate, $X_1(x),\dots,X_k(x)$ are linearly independent, spanning a
$k$-dimensional subspace $W\subseteq T_xM$. For any $i,j$,
\begin{equation}
\omega(X_i,X_j)=\PB{F_i}{F_j}=0,
\end{equation}
so $W$ is isotropic. A general fact about a $2n$-dimensional symplectic
vector space $(V,\omega)$ is that its \emph{symplectic orthogonal}
$W^{\omega}=\{v\in V:\omega(v,w)=0\ \forall w\in W\}$ satisfies
$\dim W+\dim W^\omega=2n$; isotropy of $W$ means $W\subseteq W^\omega$, so
$\dim W\le\dim W^\omega=2n-\dim W$, giving $\dim W\le n$. Hence $k\le n$.
\end{proof}

\begin{corollary}\label{cor:bound}
Of the $2n$ (general case) or $2n-1$ (autonomous case) integrals guaranteed
by Theorem~\ref{thm:universal} and Corollary~\ref{cor:autonomous}, at most
$n$ can ever be pairwise in involution, however they are chosen. Mere
existence of ``enough'' integrals, in number, therefore never by itself
implies Liouville integrability. Integrability requires exactly $n$
independent integrals \emph{in involution}: a restrictive algebraic
condition that is not satisfied generically, and is satisfied, when it is,
typically because of extra symmetry (an abelian group of symmetries acting
appropriately, or separability of the Hamilton--Jacobi equation) --- not
because of the counting argument of Section~\ref{sec:universal}.
\end{corollary}

\section{Examples}\label{sec:examples}

\begin{example}[Harmonic oscillator, $n=1$]\label{ex:sho}
Let $H=\dfrac{p^2}{2m}+\dfrac12 m\omega^2q^2$, with general solution
$q(t)=A\cos(\omega t+\varphi)$, $p(t)=-m\omega A\sin(\omega t+\varphi)$, so
$2n=2$ constants $(A,\varphi)$. Inverting,
\begin{equation}
A(q,p)=\sqrt{q^2+\Big(\frac{p}{m\omega}\Big)^2},\qquad
\varphi(q,p,t)=\operatorname{atan2}\!\Big(-\frac{p}{m\omega},\,q\Big)-\omega t
\ \ (\mathrm{mod}\ 2\pi).
\end{equation}
The amplitude $A$ (equivalently the energy $H=\tfrac12m\omega^2A^2$) is the
$2n-1=1$ time-independent integral of Corollary~\ref{cor:autonomous}; the
phase $\varphi$ is the one integral carrying explicit $t$-dependence, and
it is only defined modulo $2\pi$ and only smooth away from the equilibrium
$A=0$, illustrating the local, non-global nature of the flow-box
construction even in the simplest possible case.

For $n=1$ the two notions of Section~\ref{sec:symmetry} and
Section~\ref{sec:universal} in fact coincide: since phase space is
two-dimensional, the orbits of an autonomous flow are (generically)
one-dimensional level curves, so \emph{any} time-independent integral must
locally be a function of energy alone. The interesting distinction between
``symmetry-generated'' and ``merely guaranteed to exist'' integrals is
vacuous at $n=1$; it first becomes visible at $n\ge 2$.
\end{example}

\begin{example}[Two oscillators, $n=2$: resonance and global multivaluedness]\label{ex:pair}
Let $H=\big(\tfrac{p_1^2}{2m_1}+\tfrac12m_1\omega_1^2q_1^2\big)+\big(\tfrac{p_2^2}{2m_2}+\tfrac12m_2\omega_2^2q_2^2\big)$,
with constant frequencies $\omega_1,\omega_2$. Writing each oscillator in
action-angle form, $E_j(q_j,p_j)$ and $\theta_j(q_j,p_j,t)=\theta_j^{(0)}(q_j,p_j)+\omega_jt$
for $j=1,2$, the $2n=4$ raw constants are $E_1,E_2,\theta_1^{(0)},\theta_2^{(0)}$.
The two energies $E_1,E_2$ are already time-independent and are in
involution ($\PB{E_1}{E_2}=0$, since they depend on disjoint canonical
pairs): this accounts for $2$ of the $2n-1=3$ time-independent integrals
required by Corollary~\ref{cor:autonomous}, and, being $n=2$ independent
commuting integrals, they already establish Liouville integrability on
their own.

The third time-independent integral must be built from the two phases; the
explicit-$t$ dependence cancels in the combination
\begin{equation}
\Psi:=\omega_2\,\theta_1^{(0)}(q_1,p_1)-\omega_1\,\theta_2^{(0)}(q_2,p_2)
=\omega_2\theta_1-\omega_1\theta_2 .
\end{equation}
Two qualitatively different things can happen:
\begin{itemize}
\item If $\omega_1/\omega_2$ is rational, orbits are closed Lissajous
curves, and (a suitably normalized version of) $\Psi$ is a genuine
\emph{global}, smooth, closed-form integral --- a resonance-induced
``extra'' conserved quantity beyond $E_1,E_2$ separately, of the same
character as the classical analogue of the accidental $SU(2)$ symmetry of
the two-dimensional isotropic oscillator ($\omega_1=\omega_2$).
\item If $\omega_1/\omega_2$ is irrational, the orbit is dense on the
invariant torus $\{E_1,E_2=\text{const}\}$ (this is the classical
irrational rotation of the torus). Then $\Psi$, while still smooth and
well-defined \emph{locally}, cannot be extended to a continuous
single-valued function on the whole torus: a continuous invariant of a
dense orbit must be constant on a dense set, hence constant everywhere,
which $\Psi$ is not. The system remains Liouville integrable throughout
(the actions $E_1,E_2$ suffice), but the particular ``extra'' integral
handed to us by the general construction of Theorem~\ref{thm:universal}
simply fails to be a nice global object.
\end{itemize}
This shows explicitly that even for a fully linear, exactly solvable,
integrable system, the universal integrals of Section~\ref{sec:universal}
need not be globally smooth --- confirming the caveat after
Corollary~\ref{cor:autonomous}.
\end{example}

\begin{example}[A chaotic system: H\'enon--Heiles]\label{ex:hh}
Let $H=\tfrac12(p_x^2+p_y^2)+\tfrac12(x^2+y^2)+x^2y-\tfrac13y^3$
($n=2$), the H\'enon--Heiles Hamiltonian \cite{henonheiles}. Energy is
conserved (time-translation invariance), but no further global,
closed-form integral is known, and numerical experiments show the system
is non-integrable: above a threshold energy the accessible region of phase
space contains a chaotic sea alongside surviving invariant tori. By
Theorem~\ref{thm:universal} and Corollary~\ref{cor:autonomous}, $2n-1=3$
time-independent integrals still exist, in the strict local sense, near any
non-equilibrium point. But in the chaotic region they cannot be extended to
smooth functions on any sizeable open set, have no closed form, and are
exponentially sensitive to initial conditions (a manifestation of positive
Lyapunov exponents): two nearby initial points generically flow to values
of these ``integrals'' that diverge exponentially fast under backward
iteration of the construction. The guarantee of Theorem~\ref{thm:universal}
is not false here --- it is simply empty of practical content, in sharp
contrast with the single, robust, globally defined, closed-form Noether
integral $H$.
\end{example}

\section{Discussion}\label{sec:discussion}

The tension between statements (A) and (B) dissolves once the word
``integral'' is recognized as covering two different things:
\begin{itemize}
\item \emph{Universal integrals} (Theorem~\ref{thm:universal}): guaranteed
to exist for any system, by the bare existence-and-uniqueness theorem for
ODEs; equivalent, in information content, to already having solved the
equations of motion; generically without closed form; valid, in general,
only locally (Example~\ref{ex:pair}); with no stable functional form under
a change of Hamiltonian (perturbing $H$ at all, whether or not it breaks a
symmetry, changes the flow map $\Phi_{t,t_0}$ and hence changes every
$C_a$ completely); and, by Proposition~\ref{prop:bound}, never more than
$n$ of them can be placed in involution, so they essentially never
contribute to integrability beyond what symmetry already provides.
\item \emph{Symmetry-generated integrals} (Section~\ref{sec:symmetry}):
tied to an actual invariance of $H$; retain the exact same functional form
(e.g.\ $p_i$, or $L=q\times p$) under any deformation of $H$ that preserves
the symmetry, however drastically the rest of the dynamics is changed;
automatically in involution when they arise from an abelian symmetry group
acting appropriately, which is what makes Liouville integrability possible
in the first place.
\end{itemize}
A closely related distinction is made under the name \emph{isolating}
versus \emph{non-isolating} integrals in the stellar/galactic dynamics
literature \cite{binneytremaine}: an isolating integral (such as energy, or
angular momentum in an axisymmetric potential) genuinely confines an orbit
to a lower-dimensional invariant surface throughout the accessible region
of phase space, whereas a non-isolating integral --- guaranteed to exist by
essentially the same local argument as Theorem~\ref{thm:universal} --- has
level sets that are typically dense and provide no actual constraint on the
orbit's global behaviour. Example~\ref{ex:pair} is a clean, fully solvable
illustration of exactly this phenomenon.

The upshot is that the number $2n$ (or $2n-1$) in statement (B) is a
statement about the local invertibility of a flow, not a statement about
dynamics being simple, solvable, or governed by any symmetry. Statement (A)
is the substantive one: whether a system has symmetry-generated integrals,
and whether enough of them are in involution to achieve Liouville
integrability, is a real, non-trivial question about the specific
Hamiltonian at hand, to which the answer is generically ``no'' once
$n\ge 2$. That a bare existence theorem and a substantive dynamical
statement can be phrased using the same word, ``integral of motion,'' is
itself a small instance of a larger and more general phenomenon: results
that are transparently true once translated into pure mathematics can look
like non-trivial physical claims when left in physicists' language, and
conflating the two registers is a recurring source of apparent paradox in
classical mechanics.

\section*{Acknowledgments}
None.

\begin{thebibliography}{9}

\bibitem{landau}
L.~D.~Landau and E.~M.~Lifshitz, \emph{Mechanics}, 3rd ed., Course of
Theoretical Physics Vol.~1, Butterworth-Heinemann, 1976, \S6.

\bibitem{arnold}
V.~I.~Arnold, \emph{Mathematical Methods of Classical Mechanics}, 2nd ed.,
Springer, 1989.

\bibitem{hartman}
P.~Hartman, \emph{Ordinary Differential Equations}, 2nd ed., SIAM
Classics in Applied Mathematics, 2002.

\bibitem{binneytremaine}
J.~Binney and S.~Tremaine, \emph{Galactic Dynamics}, 2nd ed., Princeton
University Press, 2008.

\bibitem{henonheiles}
M.~H\'enon and C.~Heiles, ``The applicability of the third integral of
motion: Some numerical experiments,'' \emph{Astron.\ J.} \textbf{69}, 73
(1964).

\end{thebibliography}

\end{document}
